\subsection{The category lower bound for $\mathbb CP^2$}

\begin{lemma}[Two-object cup-product reductions]
\label{C:two-object-cup-pruning}
Let $C$ be a finite category with two objects and endomorphism monoids
$M,N$. Suppose first that all cross arrows go in one direction, and
write $S$ for that nonempty cross-biset.
\begin{enumerate}
\item If $|S|=1$, every positive-degree integral cohomology product of
      $BC$ vanishes.
\item If $|S|=2$ and one of $BM,BN$ is contractible, every product of
      two positive-degree integral classes of $BC$ is annihilated by two.
      If that vertex space is only rationally acyclic, every such
      rational product vanishes.
\item If both $BM$ and $BN$ are rationally acyclic, every positive-degree
      rational product of $BC$ vanishes, for arbitrary $S$.
\end{enumerate}
If arrows instead go in both directions and one endomorphism monoid is
trivial, $BC$ is homotopy equivalent to the classifying space of the
other endomorphism monoid. If each of the two cross-hom sets has one
element, $BC$ is contractible.
\end{lemma}
\begin{proof}
The collage homotopy-pushout description expresses $BC$ as
\[
 BM\ \cup^h_{E_S}\ BN,
\]
where $E_S$ is the two-sided double bar on the cross-biset. A singleton
biset has $E_S\simeq BM\times BN$, with the two projections, giving
$BC\simeq BM*BN$. A join of connected based spaces is homotopy
equivalent to a suspension, so its positive-degree products vanish.
If both vertex spaces are rationally acyclic, their maps to points give
a rational homology equivalence from the displayed homotopy pushout
to the unreduced suspension of $E_S$, by Mayer--Vietoris. The induced
rational cohomology ring isomorphism proves the third assertion.

For the second assertion, identify $S$ with the endpoints of the
interval $I=[0,1]$. Every self-map of the two-point set extends uniquely
to an affine map of $I$. The commuting left and right monoid actions
therefore extend to commuting interval actions, each of which is the
identity, reflection, or a constant endpoint map. The corresponding
double bar $E_I$ is homotopy equivalent over $BM\times BN$ to that
product, because $I$ is contractible. Assume $BN$ is contractible.
Collapsing that vertex and using the equivalence $E_I\to BM$ gives
\[
 BC\simeq\operatorname{hocofib}(E_S\to BM)
       \simeq\operatorname{hocofib}(E_S\to E_I).
\]
The last space is the reduced double bar with fiber
$I/\partial I=S^1$, with its basepoint bar collapsed. These comparisons
can be made as ordinary quotient comparisons: at each bar level the
endpoint inclusion is a CW cofibration, and all degeneracies are
cofibrant inclusions.

On the circle the monoid actions are identity, complex conjugation,
or the constant map to the basepoint. They all commute strictly with
$z\mapsto z^2$. This supplies a self-map $F$ of the quotient double
bar. In its reduced cellular chain complex the internal circle complex
is $\mathbb Z_\chi$ concentrated in degree one, where each monoid
element acts by $0$, $1$, or $-1$. The circle power map acts by
multiplication by two on this complex. Hence its map on the total
reduced bar chain complex is literally multiplication by two in every
degree, and $F^*$ is multiplication by two on every positive-degree
integral cohomology group. For positive classes $a,b$, multiplicativity
now gives
\[
 2ab=F^*(ab)=F^*(a)F^*(b)=4ab.
\]
Thus $2ab=0$.

If $BN$ is only rationally acyclic, the same two comparison maps are
rational homology equivalences: use Mayer--Vietoris for the first and
K\"unneth followed by the cofiber sequence for the second. Their maps
on rational cohomology are ring isomorphisms. The reduced circle-bar
argument therefore still gives vanishing of every positive rational
product. No acyclicity assertion from a finite range of homology is
used here.

For arrows in both directions and a trivial endomorphism monoid at
the second object, any pair of opposite cross arrows composes to its
identity. That object is a retract of the first. The first object is
therefore Cauchy dense, and its full-subcategory inclusion induces the
claimed equivalence on classifying spaces.

Finally suppose both cross-hom sets are singletons, with arrows
$i:1\to0$ and $r:0\to1$. Their composites $ir$ and $ri$ are zero
elements in $M,N$, respectively: multiplication by any endomorphism
leaves the singleton cross arrow unchanged. Split the idempotent $ir$
in the Karoubi category. The new object has exactly one arrow to and
from either old object, and only its identity as endomorphism; the
cross arrows identify its alternative presentation by $ri$.
It is consequently a zero object. Adding a split idempotent preserves
the classifying-space type, which proves contractibility.
\end{proof}

\begin{lemma}[Small vertex monoids in a triangular category]
\label{C:small-triangular-vertices}
No triangular two-object category with at most ten arrows has
classifying space $\mathbb CP^2$.
\end{lemma}
\begin{proof}
Write $m,n$ for its two endomorphism orders and $s\geq1$ for the
number of cross arrows. Thus $m+n+s\leq10$. Every monoid of order
at most four has finite-group homotopy type by
Theorem~\ref{E:all-six-types}, so its classifying space is rationally
acyclic. Lemma~\ref{C:two-object-cup-pruning} disposes of $s=1$,
of $s=2$ (when $\min(m,n)\leq4$), and of every case with both
$m,n\leq4$. Up to interchanging the objects, the remaining maximal
profiles are
\[
 (m,n,s)=(6,1,3),\quad(5,1,4),\quad(5,2,3).
\]
We prove statements valid for arbitrary cross-biset size in these
three endomorphism ranges.

First suppose $n=1$, so the category is the homotopy cofiber of
$E_S=B(S\rtimes M)\to BM$. If $H^2(BM;\mathbb Q)=0$, every
degree-two class of that cofiber lifts from its cohomology relative
to $BM$. The quotient by $BM$ is a suspension of $E_S$, so the
product of any two such lifted relative classes, and hence of the
original classes, is zero. The same argument applies when $BN$ is
merely rationally acyclic, by the rational cofiber comparison in
Lemma~\ref{C:two-object-cup-pruning}.

The complete order-six classification leaves only the following
monoids with nonzero rational $H^2$: the five-element sphere monoid,
its fresh-identity extension, the additional six-element sphere band,
the two regular $C_2$-band monoids of types
$S^2\vee\Sigma^2\mathbb RP^\infty$ and
$\mathbb RP^\infty/\mathbb RP^1$, and the six-element Sweeney monoid.
For all but the last, the trivial rational module has projective
dimension at most two. The sphere monoid and the additional band have
the explicit length-two resolutions already checked in
Theorem~\ref{E:all-six-types}. A fresh identity splits off an additional
rational ring factor and preserves that resolution. For either regular
group-band monoid, put $R=\mathbb QM$ and
$A=R/\mathbb QK=\mathbb QC_2$. The ideal $\mathbb QK$ is projective,
so $\operatorname{pd}_R A\leq1$. The row augmentation module $V$ is
an $A$-module and is projective over the semisimple ring $A$.
Hence $\operatorname{pd}_R V\leq1$, and the exact sequence
$0\to V\to Re\to\mathbb Q\to0$ gives
$\operatorname{pd}_R\mathbb Q\leq2$. The same statements hold for
opposites; the sphere bands and their fresh-identity extension admit
the row-column interchange, and the regular-band argument is symmetric.

For any finite right $M$-set $S$, the action category has homology
\[
 H_q(E_S;\mathbb Q)
  =\operatorname{Tor}^{\mathbb QM}_q(\mathbb QS,\mathbb Q).
\]
Thus $H_3(E_S;\mathbb Q)=0$ in these cases. Since
$H_4(BM;\mathbb Q)=0$, the cofiber also has zero fourth rational
homology and cannot be $\mathbb CP^2$.

For the Sweeney monoid use the complete resolution in
Theorem~\ref{thm:E-wedge}. Its modules $P_q$ are $P\oplus R$ for
every $q\geq2$, and its differentials satisfy
$d_{q+2}=d_q$ for $q\geq3$. The cofiber chain complex can be calculated by the algebraic
mapping cone of the natural augmentation map
\[
 \mathbb QS\otimes_R P_*\longrightarrow
 \mathbb Q\otimes_R P_*.
\]
Its chain modules and differentials are exactly two-periodic in the
range needed to compute $H_q$ for $q\geq4$. Consequently
$H_4(BC;\mathbb Q)\cong H_6(BC;\mathbb Q)$ for every finite $S$.
This excludes $\mathbb CP^2$, independently of any homology cutoff.

Finally consider $m\leq5$, $n=2$. The rational algebra of either
two-element monoid, $C_2$ or $\{1,0\}$, is semisimple. If the
larger vertex has zero rational $H^2$, the relative-cup argument above
applies. Otherwise that monoid is the five-element sphere monoid.
The double-bar edge space is the homotopy-orbit space of the
$(N,M)$-biset, and its homology is computed over
$\mathbb QN\otimes\mathbb QM^{\mathrm{op}}$. Tensoring the
length-two sphere resolution with the projective trivial
$\mathbb QN$-module gives projective dimension at most two for the
trivial module of this product algebra. The edge space again has
zero third rational homology; the two vertex spaces have zero fourth
rational homology. Mayer--Vietoris for the collage then gives
$H_4(BC;\mathbb Q)=0$. This treats every remaining profile.
\end{proof}

\begin{lemma}[Three-object star reductions]
\label{C:star-cup-pruning}
Suppose a three-object category has cross arrows only on two legs of
a star, both pointing toward its center or both pointing away from it.
If both cross-hom sets are singletons, all positive cohomology products
vanish. If one leg is a singleton with a trivial endomorphism monoid
at its leaf, and the other leaf has rationally acyclic classifying
space, all positive rational products vanish.
\end{lemma}
\begin{proof}
Write $X$ for the center's classifying space and $Y,Z$ for those of
the leaves. The nerve is the union of the two leg collages along $X$;
there are no composable pairs of cross arrows. For singleton legs
these collages are $X*Y$ and $X*Z$. Both inclusions of $X$ are
null-homotopic, by choosing a point in the respective nonempty leaf
space. Their homotopy pushout is therefore
\[
 (X*Y)\vee(X*Z)\vee\Sigma X.
\]
Each summand is a suspension, proving the first assertion.
For the second, the singleton leg ending at a point is a cone on $X$.
Collapsing this contractible collage leaves the homotopy pushout
$*\leftarrow E\to Z$, where $E$ is the other edge's double bar.
If $Z$ is rationally acyclic, its map to a point makes this pushout
rationally homology equivalent to the unreduced suspension of $E$,
with a rational cohomology ring isomorphism. This proves the claim.
\end{proof}

\begin{lemma}[Triangular categories with eleven arrows]
\label{C:triangular-eleven}
No triangular two-object category with at most eleven arrows has
classifying space $\mathbb CP^2$.
\end{lemma}
\begin{proof}
The preceding reductions leave only the new maximal profiles
\[
 (m,n,s)=(7,1,3),(6,2,3),(5,3,3).
\]
Cases with a smaller vertex of
order one or two and a larger vertex of order at most six or five,
respectively, were treated for arbitrary $s$ in
Lemma~\ref{C:small-triangular-vertices}.
For $(6,2,3)$ the same proof extends: the two-element monoid algebra
is semisimple over $\mathbb Q$, so tensoring with its projective
trivial module preserves the projective-dimension bound or the
Sweeney resolution's exact two-periodic tail. If the six-element
vertex has zero rational $H^2$, the relative-cup argument applies.

For $(7,1,3)$ we exhaust the finite three-point actions, using a small
complete structural list of target monoids. A seven-element monoid
with nonzero rational $H^2$ must have at least two minimal rows and
columns: otherwise the one-sided corner equivalence reduces it to a
finite group. Its minimal group is then trivial by the size bound,
and its minimal band is $2\times2$, $2\times3$, or $3\times2$.
The latter two cases are the band with an identity adjoined.
In the $2\times2$ case there are exactly two further nonidentity
elements. Each is determined on the band by a row transformation and
a column transformation of a two-point set. By the degree-two
balanced tensor formula, their two characters must agree for rational
$H^2$ to be nonzero. There are six such pairs of transformations.
The four products of the two extra elements may be any table entries
whose row and column translations equal the corresponding composite
translations. The program tries every such choice and checks all
$7^3$ associativity equations.

This yields $229$ normalized tables and $41$ ordinary isomorphism
classes, with exact row and column permutations and interchange of
the two extra elements used to recognize isomorphisms. Adding the
two nonsquare bands gives $43$ target monoids. For each, the program
chooses a minimal generating set and assigns every possible
transformation of a three-point set to each generator. It extends
these assignments along fixed generator words and checks all monoid
multiplication equations. Exact state permutations remove only
isomorphic actions. This exhausts $572$ right-action cases; taking
opposite monoids, already in the list, also covers the other arrow
orientation. The normalized category cup computation gives $194$
zero squares, $203$ failures of the required second Betti number,
and $173$ failures of the first Betti number. The two remaining
nonzero-square cases both have third mod-two Betti number one.
The complete generation and certificates are
\texttt{a\_CP2\_11\_triangular.py} and its JSON output.

For $(5,3,3)$, the five-element vertex must be the sphere monoid:
every other five-element monoid has zero rational $H^2$, so is
excluded by the relative-cup argument. All three-element vertex
monoids are rationally acyclic. The complete order-three census
contains seven ordinary isomorphism classes after opposites are
restored. The program \texttt{a\_CP2\_11\_bisets.py} enumerates all
three-point right actions of the sphere monoid and all left actions
of each three-element monoid, and keeps exactly the commuting pairs.
For completeness, after fixing the right action up to state
permutation it includes all conjugates of each left action, so it
does not discard a possible relative positioning of the actions.
There are $277$ resulting commuting cases. Of these, $54$ have
$b_1=0,b_2=1$, and every one has an explicit degree-three primitive
for the square. The other $223$ fail $b_1$ or $b_2$.
This completes the three profiles and the proof.
\end{proof}

\begin{proposition}[No $\mathbb CP^2$ category with at most eleven arrows]
\label{C:CP2-lower10}
\label{C:CP2-lower11}
\label{C:CP2-lower12}
Every finite category whose classifying space has the homotopy type of
$\mathbb CP^2$ has at least twelve arrows.
\end{proposition}
\begin{proof}
The one-object case is covered by the stronger monoid exclusion in
Proposition~\ref{prop:B-fourdim-lower14}. For a connected category with
$n$ arrows and $k\geq2$ objects, at least $k-1$ nonidentity arrows are
needed to connect the objects, so $k\leq(n+1)/2$. The complete
multiobject enumeration from Section~\ref{sec:C} therefore supplies
all cases through eight arrows, and all nine-arrow cases with
$3\leq k\leq5$.

For nine arrows and two objects, the additional program
\texttt{a\_categories\_cp2.cpp} uses the same complete sorted-hom-set
enumeration and associativity propagation as
\texttt{c\_categories.cpp}. It discards a hom-cardinality profile only
when it guarantees an initial or terminal object. Such a category is
contractible. At every remaining associative table, it computes the
exact normalized-nerve coboundary ranks over $\mathbb F_2$ and retains
the table precisely when $b_1=0$ and $b_2=1$. All potentially relevant
tables are therefore retained. The computation examines $8\,964\,544$
connected labelled tables after the initial/terminal-profile exclusion,
using $264\,734\,681$ propagation nodes, and retains $399\,484$ tables.
There are $776$ retained classes up to isomorphism and opposition.
As an independent check on this added filter, running it at eight arrows
produces exactly the same set of $7132$ candidate tables as the previously
completed eight-arrow census, not merely the same count.

The same complete engine at ten arrows retains respectively
$309\,960$, $6048$, and $480$ labelled tables for $k=3,4,5$.
It examines respectively $2\,591\,868$, $51\,356$, and $3900$
associative tables after the initial/terminal-profile exclusion.
These runs are complete; the two-object ten-arrow case is treated
separately below using the preceding structural lemmas.

For every retained category, \texttt{a\_category\_cups.py} constructs
the full normalized nerve in degrees through four, finds a generator
$x\in H^2(-;\mathbb F_2)$, and tests $x\smile x$ against the full image
of the degree-three coboundary. Every zero square is accompanied by an
explicit degree-three cochain primitive, checked on every normalized
four-simplex. Exact arrow relabelings and opposition alone are used to
reuse a computation; no hash invariant is treated as an isomorphism
test. The relevant counts are
\[
\begin{array}{c|c|r|r|r}
 n&k&b_1=0,\ b_2=1\text{ classes}&x^2=0&x^2\ne0\\\hline
 8&2&72&72&0\\
 8&3&23&23&0\\
 8&4&1&1&0\\
 9&2&776&773&3\\
 9&3&150&149&1\\
 9&4&8&8&0\\
 9&5&0&0&0\\
 10&3&1238&1233&5\\
 10&4&72&72&0\\
 10&5&3&3&0\\
 11&2&1&1&0\\
 11&3&2870&2842&28\\
 11&4&545&544&1\\
 11&5&22&22&0\\
 11&6&0&0&0
\end{array}
\]
All smaller multiobject cases have zero square whenever $b_1=0,b_2=1$;
their individual counts are recorded in
\texttt{a\_categories\_CP2\_small\_summary.json}.

At nine arrows, each of the three two-object residual categories has
$H^3(-;\mathbb F_2)\cong\mathbb F_2$, by the exact ranks of the full
degree-two and degree-three coboundaries. The single three-object
residual category has $H_3(-;\mathbb Z)\cong\mathbb Z/2$, and also
nonzero third mod-two homology. At ten arrows the five three-object
residual categories have third mod-two Betti numbers
$2,2,1,2,1$, respectively. Thus none of the residual categories
has the homology of $\mathbb CP^2$. Every other category in the table was excluded
by $b_1$, $b_2$, or the fact that the square of the degree-two generator
of $H^*(\mathbb CP^2;\mathbb F_2)$ is nonzero.

For two objects and ten arrows,
Lemma~\ref{C:small-triangular-vertices} removes every triangular
profile. Lemma~\ref{C:two-object-cup-pruning} removes profiles having
an endomorphism monoid of order one, by reduction to a smaller
monoid, and profiles with singleton hom sets in both directions,
by contractibility. Object exchange and category opposition allow
the remaining profiles to be normalized by
\[
 |\operatorname{End}(0)|\geq|\operatorname{End}(1)|,
 \qquad |\operatorname{Hom}(0,1)|\geq
        |\operatorname{Hom}(1,0)|.
\]
These are independent operations on the two cardinality pairs;
opposition does not change endomorphism cardinalities. The complete
program \texttt{a\_categories\_cp2\_strong.cpp} thus removes $106$
profiles by this symmetry and $39$ by the proved structural
reductions. Exactly $11$ connected profiles remain. Exhaustive
associativity propagation visits $1\,872\,751$ nodes and yields
$19\,983$ associative labelled tables. None has $b_1=0,b_2=1$.
The full output is
\texttt{a\_categories\_10\_2\_CP2\_strong.out}; its candidate
table file has only the header because the candidate set is empty.
This completes every ten-arrow case.

At eleven arrows, the program \texttt{a\_categories11\_cp2.cpp}
uses the same exhaustive propagation, with the proved reductions of
Lemmas~\ref{C:two-object-cup-pruning}, \ref{C:star-cup-pruning},
and \ref{C:triangular-eleven}. For two objects it removes $145$
profiles by object exchange and opposition and $48$ by the structural
lemmas. The $17$ remaining connected profiles produce $1\,002\,072$
associative labelled tables after $133\,327\,715$ propagation nodes.
Exactly $24$ have $b_1=0,b_2=1$; they are one isomorphism/opposition
class and have zero square. For three objects the star lemma removes
$384$ profiles, in addition to the initial/terminal-profile exclusion.
The complete counts at each object number are
\[
\begin{array}{c|r|r|r}
 k&\text{associative labelled tables}&\text{propagation nodes}
    &b_1=0,b_2=1\text{ labelled tables}\\\hline
 2&1\,002\,072&133\,327\,715&24\\
 3&28\,311\,846&626\,523\,753&1\,727\,652\\
 4&1\,030\,548&6\,178\,468&104\,016\\
 5&50\,470&132\,350&5040\\
 6&2580&2580&0
\end{array}
\]
These counts refer to the remaining profiles, so they do not silently
count structurally excluded categories as generated tables. The exact
cup-square calculation gives the eleven-arrow rows in the preceding
table. Every one of the $28$ three-object residuals has nonzero third
mod-two homology (its dimension is one, two, or three); the one
four-object residual has third mod-two Betti number one. The program
\texttt{a\_category\_residual\_ranks.py} records all normalized chain
dimensions and boundary ranks establishing these assertions. Thus
every eleven-arrow category is excluded, which proves the bound.

All complete candidate tables, normalized cocycles, explicit primitives,
and residual rank certificates are supplied in the
\texttt{a\_categories\_*\_CP2\_cups.json} files. Preliminary
nine- and ten-arrow streams marked \texttt{partial} were superseded
by the complete filtered searches and are not used for any
completeness claim.
\end{proof}
